\scanpage{19}
commutatif à isom près, on en déduit (p.ex. à l'aide de la caractérisation universelle de la prop) un \deleted{foncteur} \inserted{morphisme de topos}
\[
\mathcal L(u):\mathcal L(\alpha,\beta)\longrightarrow\mathcal L(\lambda,\mu)
\]
avec transitivité,\unsure{\ldots} rendant commutatif (à isom can. près) le diagramme
\diagram{\begin{tikzcd}[ampersand replacement=\&]
X\arrow[r,"\omega_X"]\arrow[d,"u_0"']\arrow[dr,phantom,"\overset{\sim}{\nearrow}" description]\&\mathcal L(\alpha,\beta)\arrow[d,"\mathcal L(u)=v"]\\
Y\arrow[r,"\omega_Y"']\&\mathcal L(\lambda,\mu)
\end{tikzcd}};
d'ailleurs les foncteurs $\mathcal L(u_0)_*$ et $\mathcal L(u_0)^*$ s'explicitent essentiellement comme indiqué au n°1, qu'on \deleted{\unsure{expl}} interprète \deleted{par le} en disant qu'on est dans le cas
\[
X'=X,\quad Y'=Y,\quad\beta=\mathrm{id}_X,\quad\mu=\mathrm{id}_Y,\qquad\alpha=f_X,\quad\lambda=f_Y,
\]
et en disant qu'on a un carré commutatif comme dessus\ldots\ldots

Revenons au cas général. Si $\mathcal E\inserted{=(E,F_E)}$ est un objet abélien de $\mathcal L(\alpha,\beta)$, on a des ``hom de chgt. de base''
\[
\omega_X^*(R^iv_*(\mathcal E))\longrightarrow R^iu_{0*}(\omega_X^*(\mathcal E))=R^iu_{0*}(E),
\]
et on aimerait savoir que ce sont des \uline{isoms}, sachant que c'est le cas pour $i=0$. On est obligé pratiquement d'établir que \deleted{\unsure{de}} si $\mathcal E\inserted{=(E,F_E)}$ est un faisceau abélien \uline{injectif} de $\mathcal L(\alpha,\beta)$, alors $E$ est un faisceau acyclique \deleted{\unsure{en}} (SGA4 V\quad). Il suffit pour ceci que $\alpha$ et $\beta$ commutent aux produits infinis (i.e. $\alpha_!$, $\beta_!$ existent)
\mspage{18}
